% Section IV: Experimental Setup
% All numbers sourced from PHASE_B_FINDINGS.md and experiments/out/*.json
% (record/eval files named in comments where non-obvious).

\section{Experimental Setup}
\label{sec:setup}

\subsection{Datasets}
\label{sec:setup:datasets}

We evaluate on four public aerial datasets spanning native UAV LiDAR,
UAV photogrammetry, and controlled synthetic photogrammetry. Figures
show pipeline products, not dataset ground truth; we release code,
configuration, and evaluation scripts that reproduce every number from
the original data, but no derived products of the non-commercially
licensed sources.

\subsubsection{Hessigheim 3D (H3D)}
\looseness=-1 The primary dataset~\cite{koelle2021h3d}: UAV laser scanning (Riegl
VUX-1LR, octocopter) over Hessigheim, Germany, at roughly
$800$~points/m$^2$, with per-point RGB from oblique imagery
and eleven manual semantic classes (CC~BY-NC-SA~4.0;
registration-gated). We use the three UAV epochs: March~2018
($59.4\times10^{6}$ training points plus a spatially disjoint
$14.5\times10^{6}$-point validation area), November~2018
($128.9\times10^{6}$), and March~2019 ($119.0\times10^{6}$). These
support cross-epoch replication and multi-epoch change detection
(Secs.~\ref{sec:results:seams}--\ref{sec:results:change}). H3D also
ships a March~2018 reference mesh; its 43 labeled training-split tiles
($9.0\times10^{6}$ faces after exact de-duplication of slab-boundary
faces; $132{,}622$~m$^2$) serve as the absolute geometric reference
(Sec.~\ref{sec:results:geoacc}). The mesh is the dataset authors'
reconstruction from the same campaign's imagery and LiDAR: an
independent processing chain over the same acquisition, not an
independent survey.

\subsubsection{Mill 19 Rubble}
\looseness=-1 An ungated UAV image collection over a construction/debris site near
Pittsburgh~\cite{turki2022meganerf}: 1{,}678 4K images with shipped
PixSfM-refined poses (publicly released; no formal license file, noted
for completeness). No point cloud ships with it, so we build a sparse
metric anchor ourselves, CPU-only (SIFT at $4\times$ downscale,
sequential-plus-spatial matching, COLMAP
\texttt{point\_triangulator}~\cite{schonberger2016colmap} with poses
held fixed, no bundle adjustment), yielding
$1{,}750{,}606$ colored points over a metrically coherent
$275\times318$~m site; this image-side scene carries the re-flight
planning experiment (Sec.~\ref{sec:results:reflight}).

\subsubsection{STPLS3D (synthetic subset)}
Synthetic aerial photogrammetry~\cite{chen2022stpls3d} (CC~BY-NC-SA~3.0,
form-delivered download): colored point clouds at ${\sim}0.1$~m spacing with 18
semantic classes, from simulated UAV flights over scenes ranging from
empty terrain to dense urban. Content varies while acquisition is held
constant, making it the natural controlled testbed for the capacity
study: 29
scenes at matched $N=3\times10^{6}$ points across a concentration and
prior-scale grid (261 controlled fits; Sec.~\ref{sec:disc-khat}).

\subsubsection{GauU-Scene V2 (SMBU)}
One scene from a DJI Matrice~300 / Zenmuse~L1 RGB-LiDAR
campaign~\cite{xiong2024gauuscene} (access by e-mailed license request): the
merged SMBU campus cloud of $169.1\times10^{6}$ colored points, our
largest single input, stressing the scaling claim and providing the
second-site held-out check
(Secs.~\ref{sec:results:seams}--\ref{sec:results:calibration}).

\subsection{Model Configuration}
\label{sec:setup:config}

\looseness=-1 All experiments run the tiled DP-Splat pipeline of
Sec.~\ref{sec:method} on an Apple-silicon laptop (36~GB RAM, CPU
only). Scenes are partitioned into a ground-plane grid
targeting ${\sim}10^{7}$ points per tile (H3D: $2\times3$ to $2\times7$
tiles; SMBU: 20 tiles; Mill~19: $10^{5}$-scale tiles matched to its
smaller anchor cloud), each tile dilated by a 5~m halo with ownership
down-weighting and concentration allocation as in
Sec.~\ref{sec:tiled-fit}. Tiles above $2\times10^{6}$ points are fitted
by the stochastic variational schedule (batch size $65{,}536$, three
data epochs, step size $\rho_t=(t+64)^{-0.7}$ as
in~\cite{dong2026dpsplat}); smaller tiles run full CAVI to a relative ELBO tolerance
of $10^{-6}$. (The SMBU production fit records SVI threshold
$5\times10^{5}$ and tile target $8\times10^{6}$. Its realized tiles
are highly uneven, from near-empty to $41.2\times10^{6}$ points: 15
ran the SVI schedule, the three smallest, each under $4\times10^{5}$
points, ran full CAVI, and two near-empty tiles were skipped. The
threshold difference matters only for the two tiles of $1.2$ and
$1.6\times10^{6}$ points, which ran SVI where the H3D threshold would
assign full CAVI.)

\looseness=-1 The \emph{production configuration} for all headline experiments is
$\alpha_{\mathrm{global}}=100$, truncation $T=256$ per tile, justified
by the capacity study of Sec.~\ref{sec:results:calibration}:
relative to the low-capacity reference arm ($\alpha_{\mathrm{global}}=1$,
$T=64$), production improves held-out predictive quality, resolves
the one calibration defect we observed, and activates the cross-tile
matcher, while error ranking is insensitive to the choice. Both arms
are reported wherever they differ materially; runtimes in
Sec.~\ref{sec:results:seams}.

\subsection{Evaluation Protocols}
\label{sec:setup:protocols}

\subsubsection{Held-out complement protocol}
\looseness=-1 In-sample density scores flatter any generative fit, so all headline
predictive numbers use a hold-out: each H3D epoch (and, for the
second-site check, SMBU) is refit on a
recorded 90\% subsample (fixed seed), and $200{,}000$ evaluation points
are drawn from the \emph{exact complement} of the recorded fit indices
under an independent seed: points the model has never seen, from the
same flight. All methods are scored on the identical point set.
The hold-out is sample-wise, not spatial: at ${\sim}800$
points/m$^2$ a complement point lies centimeters from training points,
so absolute scores and coverage describe quasi-interpolative
generalization within a flight; method \emph{comparisons} are
unaffected, and a spatially
blocked protocol is an open item (Sec.~\ref{sec:disc-open}). Each arm
is likewise a single fixed-seed refit scored once, so
Tables~\ref{tab:seam} and~\ref{tab:coverage} carry no paired
standard errors.
Every metric is reported on all points and on the \emph{seam subset}:
points lying inside two or more halo-dilated tiles, i.e., within one
halo width of a tile boundary (8.5\% of the March~2018 hold-out;
$n=16{,}941$).

\subsubsection{Seam baselines}
The stitched predictive is compared against the two hard-selection rules
of current practice (Sec.~\ref{sec:related-tiled}): \emph{crop}, which
truncates each tile's mixture to its owned box, and \emph{Voronoi},
which assigns each component to its nearest tile center and deletes it
elsewhere. Both are evaluated as normalized densities on the same
per-tile posteriors, isolating the seam-handling rule itself: a
controlled ablation on shared halo-down-weighted fits, not an
end-to-end reproduction of any native pipeline, which would train each
expanded cell at full weight before cropping; on the
equal-area grids used here the two rules coincide exactly (verified;
they differ on irregular partitions).

\subsubsection{Change-detection protocol}
\looseness=-1 The November~2018 and March~2019 production fits are compared on a
shared 1~m ground grid (0.5~m and 2~m in ablation), retaining cells with
at least 20 points in both epochs, under two complementary tests.
(i)~\emph{Null calibration}: cells whose dominant class is impervious in
both epochs, with dominant-class purity above $0.9$ in each
($n=17{,}592$), are physically stable ground; their
exceedance of the nominal-5\% threshold measures calibration
directly. (ii)~\emph{Detection}: among cells with per-epoch purity
above $0.6$, cells whose dominant class changes
between epochs (vegetation-excluded headline) serve as a labeled-change
proxy, scored by ROC AUC and precision/recall at the threshold; the
proxy undercounts true vertical change (same-class construction) and
includes vertical-less transitions (vehicles), so precision at the
threshold is not a false-alarm rate; the null test carries that role.
Independently labeled epochs mean dominant-label flips can
also arise from labeling inconsistency, boundary rasterization, or
leaf-off/leaf-on appearance without physical change; the proxy is
unvalidated and used only for ranking, never calibration.
Of $113{,}020$ valid cells, the $0.6$ purity filter leaves $93{,}144$
eligible; excluding cells whose dominant class in either epoch is Low
Vegetation, Shrub, or Tree leaves $43{,}221$ scored ($734$ positives).
A registration term $\sigma_{\mathrm{reg}}$ enters the level-of-detection
variance once (Sec.~\ref{sec:setup:metrics}); results are reported at
$\sigma_{\mathrm{reg}}=0$ and, as a sensitivity, at 2~cm.

\subsubsection{Re-flight protocol (leakage-proofed)}
\looseness=-1 On Mill~19 Rubble the 1{,}678-image flight is split by acquisition
order: the first 80\% forms the survey, the trailing 20\% block (336
images) the held-back re-flight pool, emulating an unfinished mission
over contiguous new ground. Triangulated points are assigned to survey
or pool by which image set observes them in at least two views
($1{,}480{,}265$ survey; $266{,}168$ pool; ambiguous points discarded);
an interleaved every-fifth-image split is the null control. After a
robust bounding-box crop of the anchor (the per-axis $0.5$--$99.5$
percentile box of the survey cloud, applied to both sets; it drops
$42{,}785$ survey and
$14{,}148$ pool outliers), a fixed 25\% of the cropped pool
($63{,}005$ points) is reserved as the evaluation set \emph{before}
any granting, identically across arms, so no evaluation point enters
any refit. Tiles are a ground-plane partition of the full cropped
site's bounding box: $2\times3$ in the 6-tile scenario, $3\times4$ in
the 12-tile scenario. Planner arms
select the top quarter of tiles by acquisition score
(nearest-integer rounding: 2 of 6 tiles, 3 of 12) and
receive the remaining pool points there; controls are random
tile subsets of the same size (5--8 seeded replicates) and a per-tile
\emph{oracle} ranked by the pool points' true held-out error
contribution under the initial model; improvement is the change in mean
held-out NLPD, baseline minus arm, on the fixed evaluation set.

\subsection{Metrics}
\label{sec:setup:metrics}

\subsubsection{Held-out log-predictive}
% Mass numbers (0.980, SE 0.002; per-tile-renormalized 5.62):
% experiments/out/mass_oracle_mar18.json.
Mean log-density per point (nats, i.e., natural-log units) of the
six-dimensional
(spatial~$\times$~color) stitched predictive on held-out points; higher
is better; reported for all methods on both splits. Density comparisons
require normalized densities, so the stitched mass is oracle-enforced:
tile sub-mixtures carry their components' \emph{global} weights,
un-renormalized within tiles (a straddling component contributes
$\pi_k \sum_t g_t = \pi_k$), and importance-sampling integration of the
production March~2018 model measures total mass $0.980$ (SE $0.002$;
the deficit is Student-$t$ tail mass outside the dilated footprint).
Renormalizing per tile instead inflates the mass to
${\approx}$\,\#tiles ($5.62$ measured), a regression-guarded failure
mode. The crop/Voronoi baselines renormalize to exactly unit mass, so
the stitched rule carries a $0.020$-nat ($-\log 0.980$) tail-mass
handicap into every overall comparison of Table~\ref{tab:seam}. Set
against the production deficits there, the handicap roughly cancels
the March~2018 gap ($-0.017$) but not the other two
($-0.031$/$-0.033$), which would shrink to roughly
$-0.011$/$-0.013$ if their unmeasured tail mass matches
March~2018's; the reference-arm seam-band gaps are more than an order
of magnitude larger and are not explained by it.

\subsubsection{Calibration coverage}
For each held-out point, responsibilities over all (tile, component)
pairs are formed under the stitched predictive (conditioning on the
full point, position and color), and the point is tested against each
claiming component's central equal-tailed Student-$t$ interval at
nominal 68/90/95\% along each spatial axis; reported coverage is the
responsibility-weighted mean of these indicators, averaged over axes,
on both splits. This is a component-conditional calibration check,
asking whether a point falls in the central interval of the components
that claim it, not the mixture-marginal CDF; we make no theoretical
seam-coverage claim.

\subsubsection{Sparsification: AUSE and AURG}
\looseness=-1 Following the sparsification-curve
methodology~\cite{ilg2018uncertainty,poggi2020uncertainty}, $200{,}000$
points are sampled \emph{exactly} from the stitched spatial predictive
(rejection sampling with the gates as acceptance probability), and each sample is
scored by the reference-based error convention of Huang and
Qin~\cite{huang2025uncertainty}: $e$ is the orthogonal distance to the
least-squares plane through the sample's six nearest neighbors in the
full epoch cloud. Samples are removed by decreasing uncertainty $u$ in
2\% steps up to 98\%, recording the mean $|e|$ of the retained set,
normalized by the full-set mean. AUSE (lower is better) is the area
between the uncertainty-ordered and error-oracle curves; AURG (higher
is better) that between the random-removal and uncertainty-ordered
curves. Three $u$ variants are reported: per-point NLPD under the
stitched mixture (the headline ranker); the generating component's
marginal predictive variances pooled by root mean square, written
$\sqrt{\overline{\mathrm{Var}}}$ (the location-queryable
alternative); and, as a model-free control, the negative of the local
point density $6/(\tfrac{4}{3}\pi r_6^{3})$, with $r_6$ the distance to the
sample's sixth reference neighbor: the same neighborhood the error
plane is fitted to. Density-stratified AUSE is reported alongside as
the honest confound decomposition, and the AUSE machinery
itself is verified in Sec.~\ref{sec:results:sparsification}. Unlike
the predict-then-observe setting
of~\cite{ilg2018uncertainty,poggi2020uncertainty}, the scored points
are the model's own samples ($u$ and $e$ are functionals of the
same point and model), so the ranking is model criticism, a
quality-assurance ordering of reconstructed mass against the reference
geometry, not a before-measurement error map
(Sec.~\ref{sec:results:sparsification}).

\subsubsection{Geometric accuracy and completeness}
\looseness=-1 Against the March~2018 reference mesh
(Sec.~\ref{sec:setup:datasets}), both directions of the standard
accuracy/completeness pair are scored in meters.
\emph{Accuracy} (model${\to}$mesh): $200{,}000$ samples drawn from
each seam rule's spatial predictive (stitched by the gate-rejection
sampler above; crop/Voronoi ancestrally from their normalized
mixtures), each scored by exact unsigned point-to-mesh distance (a
certified nearest-triangle search, oracle-tested against brute force);
the 1.2--1.5\% of samples falling outside the mesh tiles' footprint
are excluded. \emph{Completeness} (mesh${\to}$model): $200{,}000$
area-weighted mesh-surface samples, restricted to 1~m ground cells
holding at least five points of the fit's own input cloud (70.3\%
kept; the mesh extends beyond the surveyed area), each scored by
distance to the nearest of $5\times10^{5}$ dense model samples, a
discretization that overestimates distance to the model by about the
dense-bank spacing (median $0.326$~m), the resolution floor of the
completeness fractions. Both directions are additionally stratified
into terciles of the stitched model's own per-sample spatial
predictive NLPD.

\subsubsection{Level of detection (LoD95)}
For change detection we adopt the significance-testing convention of the
M3C2 lineage~\cite{lague2013m3c2,winiwarter2021m3c2ep}: a per-cell threshold
\begin{equation}
\mathrm{LoD}_{95}
 = 1.96\sqrt{\operatorname{var}^{A}_{z} + \operatorname{var}^{B}_{z}
   + \sigma_{\mathrm{reg}}^{2}},
\label{eq:lod95}
\end{equation}
where $\operatorname{var}^{E}_{z}$ is the vertical variance of the
responsibility-weighted \emph{posterior mean location} in epoch $E$, derived
in closed form from component $k$'s spatial NIW posterior as
$\Sigma_{\mu,k}=\Psi_{k,s}/\bigl(\kappa_{k,s}(\nu_{k,s}-D_s-1)\bigr)$. Using the
mean-location covariance rather than the full predictive covariance
avoids double-counting surface roughness in the threshold (the
pitfall M3C2-EP was designed to fix~\cite{winiwarter2021m3c2ep}), and
we verify the two covariances are distinct by Monte Carlo. Concretely,
each cell is scored per epoch at its mean observed ground point
$q = (c_x, c_y, \bar z_E)$: responsibilities
$w_{tk} \propto g_t(q)\, \mathbb{E}_q[\pi_k]\,
\operatorname{St}_3(q)$ over (tile, component) pairs (global
post-merge weights, as in \eqref{eq:stitch}) give
$\mu^E_z = \sum_{tk} w_{tk}\, m_{k,z}$, with
$\operatorname{var}^E_z$ by the law of total variance (weighted
$[\Sigma_{\mu,k}]_{zz}$ within components plus the spread of $m_{k,z}$
between them); $d = \mu^A_z - \mu^B_z$ is a vertical difference on a
fixed grid (DoD-style, conflating horizontal displacement on slopes
and facades with elevation change, unlike M3C2's normal-direction
projection), and change is flagged where $|d| > \mathrm{LoD}_{95}$.
The \emph{classical} baseline of Sec.~\ref{sec:results:change} is the
propagated-error DoD convention on identical cells, masks, and proxy
sets: $d = \bar z_A - \bar z_B$ from the raw per-cell points, and
\eqref{eq:lod95} with the empirical standard errors $s^2_E/n_E$ of
those means (within-cell sample variance over the $n_E$ points) in
place of the posterior mean-location variances; the arms differ
only in the variance source.
